\documentclass[12pt,a4paper]{article}
\usepackage[utf8]{inputenc}
\usepackage[T1]{fontenc}
\usepackage[french]{babel}
\usepackage{geometry}
\usepackage{titlesec}
\usepackage{parskip}
\usepackage{hyperref}
\usepackage{xcolor}
\usepackage{enumitem}
\usepackage{fancyhdr}
\usepackage{booktabs}
\usepackage{array}

\geometry{margin=2.5cm}

\definecolor{codeblue}{RGB}{0,70,150}
\definecolor{codered}{RGB}{180,0,0}
\definecolor{codegray}{RGB}{100,100,100}

\pagestyle{fancy}
\fancyhf{}
\rhead{Application de statistiques sportives}
\lhead{Module \texttt{src/Analysis}}
\cfoot{\thepage}

\titleformat{\section}{\large\bfseries\color{codeblue}}{}{0em}{}[\titlerule]
\titleformat{\subsection}{\normalsize\bfseries\color{codered}}{}{0em}{}
\titleformat{\subsubsection}{\normalsize\itshape}{}{0em}{}

\newcommand{\fonction}[1]{\textbf{\texttt{#1}}}
\newcommand{\module}[1]{\textit{\texttt{#1}}}

\begin{document}

\begin{titlepage}
    \centering
    \vspace*{3cm}
    {\huge\bfseries Documentation technique\\[0.5em]}
    {\Large\bfseries Module \texttt{src/Analysis}\\[2em]}
    {\large Application de statistiques sportives\\[1em]}
    {\normalsize Projet 1A — 2026\\[3em]}
    \hrule
    \vspace{1em}
    {\small Ce document décrit l'ensemble des classes et fonctions du dossier
    \texttt{src/Analysis}, responsable du calcul des statistiques et de
    l'affichage des résultats pour chaque sport.}
\end{titlepage}

\tableofcontents
\newpage

% ============================================================
\section{Architecture générale}

Le dossier \texttt{src/Analysis} organise chaque sport dans un module dédié.
Tous les modules suivent le même pattern structurel~:

\begin{enumerate}
    \item Des variables globales \texttt{\_loader}, \texttt{\_players},
    \texttt{\_matches} (etc.) initialisées à \texttt{None}.
    \item Une fonction \texttt{\_load()} qui charge les données
    \textbf{une seule fois} (cache dit \textit{lazy loading}).
    \item Des fonctions d'analyse qui appellent \texttt{\_load()} en première
    instruction.
    \item Une fonction \texttt{get\_agenda\_data()} qui formate les matchs dans
    un format standard pour l'agenda.
\end{enumerate}

Le fichier \texttt{générique.py} est le \textbf{socle partagé} entre tous
les modules~: il fournit les fonctions d'affichage (bracket, classement),
de formatage (roster, fiche joueur) et d'agrégation (agenda).

% ============================================================
\section{\texttt{générique.py} — Le moteur commun}

\textbf{Fichier~:} \texttt{src/Analysis/générique.py}

Ce fichier ne traite aucun sport en particulier. Il définit les briques
réutilisables utilisées par l'ensemble des modules sportifs.

\subsection{Classe \texttt{ResultatMatch}}

Classe utilitaire qui encapsule le résultat d'un match à partir de deux scores.

\subsubsection{Constructeur \texttt{\_\_init\_\_(equipe1, equipe2, score1, score2)}}
Calcule automatiquement les attributs~:
\begin{itemize}
    \item \texttt{joue} (bool) — \texttt{True} si les deux scores sont disponibles
    et non \texttt{NaN}. La détection du NaN utilise la technique Python
    \texttt{score != score} (un float NaN n'est jamais égal à lui-même).
    \item \texttt{gagnant} — nom de l'équipe gagnante, ou \texttt{None} si le
    match n'est pas joué ou en cas d'égalité.
    \item \texttt{perdant} — nom de l'équipe perdante, ou \texttt{None}.
\end{itemize}

\subsubsection{Méthode \texttt{label\_equipe(equipe)}}
Retourne la chaîne \texttt{"score1-score2"} si le match est joué, sinon
\texttt{"???"}.

\subsection{\texttt{formater\_roster(df\_joueurs, col\_nom, ...)}}

Formate le roster d'une équipe selon la spécification standard~:
\texttt{Rôle | Nom complet | Pseudo* | Nationalité | Date de naissance}.

\begin{itemize}
    \item Les coachs apparaissent en tête de liste, puis les joueurs.
    \item La colonne \texttt{Pseudo} n'est incluse que si \texttt{est\_esport=True}.
    \item Les colonnes absentes des données sont affichées comme \texttt{"N/A"}.
    \item Les dates sont formatées en \texttt{JJ/MM/AAAA}.
\end{itemize}

\subsection{\texttt{\_agreger\_deux\_manches(df, ...)}}

Fonction privée. Pour les compétitions aller/retour (ex~: Champions League),
additionne les scores des deux matchs entre les mêmes équipes en un seul
résultat. Utilise \texttt{tuple(sorted([eq1, eq2]))} pour identifier une paire
indépendamment du sens domicile/extérieur.

\subsection{\texttt{afficher\_bracket(df, col\_equipe1, col\_equipe2, ...)}}

Affiche un bracket d'élimination directe en \textbf{ASCII colorisé} dans le
terminal. Construit une grille 2D de chaînes, puis des colonnes de connecteurs
entre les rounds. Le chemin du gagnant est affiché en vert via les codes ANSI.
Le paramètre \texttt{deux\_manches} déclenche l'agrégation aller/retour avant
affichage.

\subsection{\texttt{\_calculer\_classement(df, ...)}}

Fonction privée. Calcule un classement par points à partir d'un DataFrame de
matchs.
\begin{itemize}
    \item Règles~: 3~pts victoire, 1~pt nul, 0~pt défaite.
    \item Tri~: Points décroissants, puis différence de buts/scores décroissante.
    \item Rang~: classement standard (ex-æquo = même rang, le suivant saute
    des places).
\end{itemize}
Retourne les colonnes \texttt{Rang | Équipe | MJ | V | N | D | Pts | Diff |
Winrate}.

\subsection{\texttt{\_afficher\_table\_classement(df, top\_qualifies, titre)}}

Fonction privée. Affiche un tableau de classement avec une ligne séparatrice
optionnelle après les \texttt{top\_qualifies} premières équipes
(ex~: \texttt{─── TOP 2 ───}).

\subsection{\texttt{afficher\_classement(df, col\_equipe1, col\_equipe2, ...)}}

Orchestre \texttt{\_calculer\_classement} et \texttt{\_afficher\_table\_classement}.
Si \texttt{col\_groupe} est fourni, itère sur chaque groupe et affiche un
tableau de classement par groupe.

\subsection{\texttt{fiche\_joueur(df\_joueurs, col\_nom, nom\_joueur, ...)}}

Retourne la fiche complète d'un joueur sous forme \texttt{Statistique | Valeur}.
\begin{itemize}
    \item Recherche insensible à la casse, accepte les noms partiels.
    \item Lève \texttt{ValueError} si aucun résultat.
    \item Lève \texttt{ValueError} si plusieurs joueurs correspondent
    (liste les correspondances pour affiner la recherche).
    \item Les colonnes sont renommées selon le dictionnaire \texttt{col\_labels}.
\end{itemize}

\subsection{\texttt{lister\_joueurs(df\_joueurs, col\_nom, ...)}}

Retourne la liste simple de tous les joueurs d'un sport, triée par nom.
Colonnes~: \texttt{Nom} (et \texttt{Équipe} si la colonne est fournie).

\subsection{\texttt{agenda\_recents(sources, n=10)}}

Agrège les DataFrames de plusieurs sports au format standard
\texttt{Sport | Date | Équipe 1 | Équipe 2 | Score 1 | Score 2}, trie par
date décroissante, et retourne les \texttt{n} matchs les plus récents.

% ============================================================
\section{\texttt{badminton.py}}

\textbf{Fichier~:} \texttt{src/Analysis/badminton.py}

\subsection{\texttt{\_load()}}
Cache lazy~: instancie \texttt{BadmintonLoader} et charge joueurs + matchs
une seule fois dans les variables globales du module.

\subsection{\texttt{\_compter\_jeux(row)}}
Fonction privée. Parse les colonnes \texttt{game\_1\_score},
\texttt{game\_2\_score}, \texttt{game\_3\_score} au format \texttt{"21-15"}
et retourne un tuple \texttt{(wins\_joueur1, wins\_joueur2)}. Gère les sets
non joués (NaN ou chaîne vide).

\subsection{\texttt{fiche\_joueur\_badminton(nom)}}
Délègue à \texttt{générique.fiche\_joueur}. Affiche nom, pays, continent.

\subsection{\texttt{bilan\_joueur\_badminton(nom)}}
Calcule victoires et défaites en cherchant le joueur comme \texttt{player\_1}
ET comme \texttt{player\_2} dans les matchs. Utilise \texttt{\_compter\_jeux}
pour déterminer le gagnant. Retourne un DataFrame
\texttt{Statistique | Nom du joueur}.

\subsection{\texttt{bracket(tournament\_name)}}
Filtre les matchs d'un tournoi sur les rounds bracket (\texttt{Round of 32}
jusqu'à \texttt{Final}), calcule les scores avec \texttt{\_compter\_jeux},
puis délègue à \texttt{générique.afficher\_bracket}.

\subsection{\texttt{get\_agenda\_data()}}
Retourne les matchs au format standard agenda, avec le nombre de jeux gagnés
comme score.

% ============================================================
\section{\texttt{basketball.py}}

\textbf{Fichier~:} \texttt{src/Analysis/basketball.py}

\subsection{\texttt{\_load()}}
Cache lazy~: charge joueurs, équipes et matchs NBA.

\subsection{\texttt{top\_equipes\_offensives(n=10)}}
Fusionne les colonnes \texttt{pts\_home} et \texttt{pts\_away} en un format
neutre \texttt{team\_id / pts}, calcule la moyenne de points par match par
équipe, trie par ordre décroissant.

\subsection{\texttt{stats\_equipe(team\_name)}}
Recherche l'équipe par nom complet, abréviation ou surnom. Fusionne les matchs
à domicile et à l'extérieur (en supprimant les suffixes \texttt{\_home} et
\texttt{\_away}). Calcule la moyenne de 12 statistiques (points, rebonds,
passes, tirs, lancers francs, interceptions, contres, pertes de balle).

\subsection{\texttt{roster\_equipe(team\_name)}}
Appelle \texttt{BasketballLoader.get\_team\_roster()} puis délègue à
\texttt{générique.formater\_roster}.

\subsection{\texttt{classement\_defensif(n=10)}}
Calcule les points encaissés, blocks et interceptions moyens par équipe.
Trie par points encaissés croissants (moins on en prend, meilleure est la
défense).

\subsection{\texttt{fiche\_joueur\_basketball(nom)}}
Délègue à \texttt{générique.fiche\_joueur} avec les labels en français.

\subsection{\texttt{liste\_joueurs(equipe=None)}}
Liste tous les joueurs NBA, ou uniquement ceux d'une équipe si un nom est
précisé.

\subsection{\texttt{get\_agenda\_data()}}
Mappe les identifiants d'équipes (\texttt{team\_id}) vers les noms complets,
retourne le format standard agenda.

\subsection{\texttt{classement\_points(saison=None, top\_qualifies=None)}}
Filtre sur \texttt{Regular Season}, mappe les identifiants d'équipes vers
les noms, délègue à \texttt{générique.afficher\_classement}.

% ============================================================
\section{\texttt{chess.py}}

\textbf{Fichier~:} \texttt{src/Analysis/chess.py}

\subsection{\texttt{\_load()}}
Cache lazy~: charge joueurs et matchs d'échecs.

\subsection{\texttt{classement\_elo(mode="standard", n=20)}}
Accepte les modes \texttt{"standard"}, \texttt{"rapid"} ou \texttt{"blitz"}.
Mappe vers la colonne Elo correspondante, supprime les NaN, trie et retourne
les N premiers joueurs avec leur fédération et titre FIDE.

\subsection{\texttt{fiche\_joueur\_chess(nom)}}
Délègue à \texttt{générique.fiche\_joueur} avec labels~: Nom, Fédération,
Titre FIDE, Elos standard/rapid/blitz.

\subsection{\texttt{bilan\_joueur(player\_name)}}
Cherche le joueur comme \texttt{player\_1} et comme \texttt{player\_2}.
Utilise \texttt{pd.to\_numeric(..., errors="coerce")} pour convertir les
scores (pouvant contenir \texttt{"Bye"}). Un score de~\texttt{1} = victoire,
\texttt{0.5} = nul, \texttt{0} = défaite. Retourne aussi le score total
(victoires + 0,5 × nuls), fidèle aux règles des échecs.

\subsection{\texttt{stats\_par\_titre()}}
Groupe les joueurs par titre FIDE (\texttt{GM}, \texttt{IM}, etc.) et calcule
le nombre de joueurs ainsi que l'Elo moyen dans chaque catégorie
(standard, rapid, blitz).

% ============================================================
\section{\texttt{cs2.py}}

\textbf{Fichier~:} \texttt{src/Analysis/cs2.py}

\subsection{\texttt{\_load()}}
Cache lazy~: charge joueurs, coachs, équipes et matchs CS2.

\subsection{\texttt{classement()}}
Calcule victoires et matchs joués par équipe sur l'ensemble de la compétition.
Détermine le gagnant de chaque match par comparaison directe des scores.

\subsection{\texttt{stats\_equipe(team\_name)}}
Calcule matchs joués, victoires, défaites, et \textbf{maps gagnées et perdues}
(cumul des scores). La distinction matchs/maps est propre au CS2~: un score
de \texttt{2-1} signifie 2 maps à 1.

\subsection{\texttt{roster\_equipe(team\_name)}}
Combine joueurs et coachs filtrés par équipe, délègue à
\texttt{générique.formater\_roster} avec \texttt{est\_esport=True}
(inclut la colonne Pseudo).

\subsection{\texttt{fiche\_joueur\_cs2(nom)}}
Délègue à \texttt{générique.fiche\_joueur}.

\subsection{\texttt{liste\_joueurs(equipe=None)}}
Liste tous les joueurs CS2, filtrables par équipe.

\subsection{\texttt{get\_agenda\_data()}}
Retourne le format standard agenda avec les scores en nombre de maps.

\subsection{\texttt{bracket()}}
Filtre les matchs de stage \texttt{"PlayOffs"} et délègue à
\texttt{générique.afficher\_bracket} avec \texttt{deux\_manches=False}
(pas d'aller/retour en CS2).

\subsection{\texttt{classement\_stages(top\_qualifies=None)}}
Affiche un classement par phase de groupes (stages 1, 2, 3). Crée une colonne
\texttt{stage\_label} et exclut les PlayOffs avant de déléguer à
\texttt{générique.afficher\_classement}.

% ============================================================
\section{\texttt{football\_cl.py}}

\textbf{Fichier~:} \texttt{src/Analysis/football\_cl.py}

\subsection{\texttt{\_load()}}
Cache lazy~: charge joueurs, équipes et matchs Champions League.

\subsection{\texttt{meilleurs\_buteurs(n=10)}}
Trie les joueurs par buts et calcule le pourcentage de tirs cadrés
(\texttt{on\_target / total\_attempts}).

\subsection{\texttt{meilleurs\_passeurs(n=10)}}
Trie les joueurs par nombre de passes décisives.

\subsection{\texttt{stats\_equipe(team\_name)}}
Agrège les statistiques de tous les joueurs d'un club (somme totale) en
buts, passes décisives, minutes jouées, cartons jaunes et cartons rouges.

\subsection{\texttt{stats\_joueur(player\_name)}}
Retourne toutes les statistiques individuelles d'un joueur de la Champions
League (buts, passes, dribbles, tacles, cartons, passes tentées/réussies,
minutes).

\subsection{\texttt{stats\_gardiens(n=10)}}
Filtre les joueurs dont la position contient \texttt{Goalkeeper},
\texttt{GK} ou \texttt{Gardien}, trie par nombre d'arrêts. Retourne arrêts,
buts encaissés, clean sheets, pénaltys arrêtés.

\subsection{\texttt{fiche\_joueur\_fcl(nom)}}
Délègue à \texttt{générique.fiche\_joueur} avec un dictionnaire de labels
très détaillé (20 colonnes).

\subsection{\texttt{liste\_joueurs(club=None)}}
Liste tous les joueurs de la CL, filtrables par club.

\subsection{\texttt{classement\_groupes(top\_qualifies=2)}}
Filtre les matchs de phase \texttt{"group"} et délègue à
\texttt{générique.afficher\_classement} avec \texttt{col\_groupe="group"}.

\subsection{\texttt{get\_agenda\_data()}}
Retourne le format standard agenda.

\subsection{\texttt{bracket()}}
Filtre les matchs de phase \texttt{"knockout"} et délègue à
\texttt{générique.afficher\_bracket} avec \texttt{deux\_manches=True}
(aller/retour en Champions League).

% ============================================================
\section{\texttt{lol.py}}

\textbf{Fichier~:} \texttt{src/Analysis/lol.py}

\subsection{\texttt{\_load()}}
Cache lazy~: charge joueurs, coachs, équipes et matchs LoL.

\subsection{\texttt{stats\_equipe(team\_name)}}
Cherche l'équipe dans les colonnes \texttt{team\_blue} et \texttt{team\_red},
fusionne les matchs en renommant les suffixes \texttt{\_team\_blue} et
\texttt{\_team\_red}. Calcule la moyenne de kills, assists, deaths, gold,
turrets, dragons, barons, et la \textbf{durée moyenne en minutes}.

\subsection{\texttt{champions\_picks\_bans(n=10)}}
Identifie toutes les colonnes commençant par \texttt{pick\_} et \texttt{ban\_},
aplatit leurs valeurs, compte les occurrences de chaque champion. Retourne
un top~N par picks et bans.

\subsection{\texttt{roster\_equipe(team\_name)}}
Combine joueurs et coachs, délègue à \texttt{générique.formater\_roster}
avec \texttt{est\_esport=True}.

\subsection{\texttt{fiche\_joueur\_lol(nom)}}
Délègue à \texttt{générique.fiche\_joueur}.

\subsection{\texttt{liste\_joueurs(equipe=None)}}
Liste les joueurs LoL, filtrables par équipe.

\subsection{\texttt{get\_agenda\_data()}}
Le score est \texttt{1} pour le gagnant, \texttt{0} pour le perdant.

\subsection{\texttt{classement\_points(top\_qualifies=None)}}
Crée des colonnes \texttt{score\_blue} et \texttt{score\_red} à partir de la
colonne \texttt{winner}, puis délègue à \texttt{générique.afficher\_classement}.

% ============================================================
\section{\texttt{starcraft2.py}}

\textbf{Fichier~:} \texttt{src/Analysis/starcraft2.py}

\subsection{\texttt{\_load()}}
Cache lazy~: charge joueurs et matchs SC2.

\subsection{\texttt{fiche\_joueur\_sc2(nom)}}
Délègue à \texttt{générique.fiche\_joueur}. Inclut la colonne \texttt{race}
(Terran / Zerg / Protoss), spécifique à StarCraft~II.

\subsection{\texttt{bilan\_joueur\_sc2(nom)}}
Cherche le joueur comme \texttt{player\_1} et comme \texttt{player\_2}.
Compare directement \texttt{score\_player\_1 > score\_player\_2} pour
déterminer le gagnant.

\subsection{\texttt{get\_agenda\_data()}}
Retourne le format standard agenda avec les scores en nombre de maps.

% ============================================================
\section{\texttt{tennis.py}}

\textbf{Fichier~:} \texttt{src/Analysis/tennis.py}

\subsection{\texttt{\_load()}}
Cache lazy~: charge les quatre tables ATP et WTA.

\subsection{\texttt{\_matches(circuit)} / \texttt{\_players(circuit)}}
Fonctions helpers privées qui retournent la bonne table selon
\texttt{"ATP"} ou \texttt{"WTA"}.

\subsection{\texttt{classement\_victoires(circuit="ATP", n=20)}}
Compte victoires (\texttt{winner\_id}) et défaites (\texttt{loser\_id})
par identifiant joueur, joint les noms complets, trie et retourne les
N meilleurs.

\subsection{\texttt{stats\_joueur(player\_name, circuit="ATP")}}
Retourne victoires, défaites, pourcentage de victoires, nombre de tournois,
et la durée moyenne d'un match si la colonne \texttt{minutes} est disponible.

\subsection{\texttt{fiche\_joueur\_tennis(nom, circuit="ATP")}}
Délègue à \texttt{générique.fiche\_joueur}.

\subsection{\texttt{liste\_joueurs(circuit="ATP")}}
Liste tous les joueurs du circuit sélectionné.

\subsection{\texttt{\_compter\_sets(score)}}
Fonction privée. Parse une chaîne de score tennis comme
\texttt{"7-6(5) 6-4"}. Utilise \texttt{re.sub(r'\textbackslash(\textbackslash d+\textbackslash)', '', ...)}
pour supprimer les scores de tie-break avant de compter les sets.

\subsection{\texttt{\_standardiser\_tennis(matches, players, label)}}
Fonction privée. Mappe les \texttt{winner\_id} et \texttt{loser\_id} vers les
noms complets, calcule les sets via \texttt{\_compter\_sets}, retourne le
format agenda standard.

\subsection{\texttt{get\_agenda\_data()}}
Concatène les matchs ATP et WTA standardisés. Le score correspond aux
sets gagnés (gagnant / perdant).

\subsection{\texttt{bracket(tourney\_name, circuit="ATP")}}
Filtre les matchs d'un tournoi, mappe les identifiants en noms, calcule les
sets, filtre sur les rounds bracket (\texttt{R128} jusqu'à \texttt{F}),
délègue à \texttt{générique.afficher\_bracket}.

% ============================================================
\section{\texttt{volleyball.py}}

\textbf{Fichier~:} \texttt{src/Analysis/volleyball.py}

\subsection{\texttt{\_load()}}
Cache lazy~: charge les sept tables de volleyball.

\subsection{\texttt{\_matches(genre)} / \texttt{\_players(genre)} / \texttt{\_country\_name(code)}}
Helpers privés pour sélectionner la bonne table par genre
(\texttt{"hommes"} ou \texttt{"femmes"}), et pour résoudre un code IOC
en nom de pays complet.

\subsection{\texttt{classement(genre="hommes")}}
Calcule victoires, défaites et \textbf{sets gagnés} (critère de départage
en volleyball). Enrichit le résultat avec les noms complets de pays via
la table \texttt{\_countries}.

\subsection{\texttt{bilan\_equipe(team\_code, genre="hommes")}}
Calcule le bilan d'une équipe identifiée par un code IOC (ex~: \texttt{"FRA"}),
avec sets gagnés et sets perdus.

\subsection{\texttt{roster\_equipe(team\_code, genre="hommes")}}
Filtre par code IOC exact (pas de recherche partielle, les codes sont
standardisés), combine joueurs et coachs, délègue à
\texttt{générique.formater\_roster}.

\subsection{\texttt{fiche\_joueur\_volleyball(nom, genre="hommes")}}
Délègue à \texttt{générique.fiche\_joueur}. Inclut \texttt{height},
\texttt{birth\_place} et \texttt{nickname}.

\subsection{\texttt{liste\_joueurs(genre="hommes", equipe=None)}}
Liste les joueurs d'un genre, filtrables par code pays.

\subsection{\texttt{get\_agenda\_data()}}
Mappe les codes IOC vers les noms complets de pays. Concatène les matchs
hommes et femmes.

\subsection{\texttt{classement\_groupes(genre="hommes", top\_qualifies=2)}}
Extrait la lettre de poule depuis la colonne \texttt{stage} via le regex
\texttt{Pool ([A-Z])}, puis délègue à \texttt{générique.afficher\_classement}.

% ============================================================
\section{Synthèse des patterns communs}

\begin{center}
\begin{tabular}{>{\ttfamily}p{4.5cm} p{9cm}}
\toprule
\textbf{Pattern} & \textbf{Description} \\
\midrule
Lazy loading & Les données ne sont chargées qu'au premier appel de \_load(),
puis mises en cache dans les variables globales du module \\
get\_agenda\_data() & Chaque module expose cette fonction au format standard,
permettant l'agrégation multi-sport via générique.agenda\_recents() \\
Délégation à générique & Les fonctions d'affichage (bracket, classement, fiche,
roster) sont centralisées dans générique.py et appelées par tous les modules \\
Recherche partielle & Toutes les recherches par nom acceptent des noms partiels
et sont insensibles à la casse \\
ValueError explicite & Si aucun résultat ou si la recherche est ambiguë,
une ValueError est levée avec un message lisible \\
\bottomrule
\end{tabular}
\end{center}

\end{document}
